%
% Curriculum Vitæ/Résumé
% Nicolas Dubois <nicolas.c.dubois@gmail.com>
% http://nicolasdubois.com • @wo0dyn
%
% KOMA-Script: scrartcl
%
% This work is licensed under the Creative Commons Attribution-NonCommercial 4.0 International
% License.
% To view a copy of this license, visit http://creativecommons.org/licenses/by-nc/4.0/.
%

\documentclass[a4paper, 11pt]{scrartcl}

% Packages:
\usepackage[american, francais]{babel}
\usepackage[utf8]{inputenc}
\usepackage[pdftex]{xcolor}

\usepackage{fancyhdr}
\usepackage{graphicx}
\usepackage{url}
\usepackage{xspace}

\usepackage{packages/vendor/aeguill}
\usepackage{packages/vendor/vmargin}

\urlstyle{sf}

\usepackage[pdftex, colorlinks,%
	pdfpagelayout=SinglePage,%
	pdftitle={Curriculum Vitæ — Résumé},%
	pdfauthor={Nicolas Dubois, @wo0dyn},%
	pagebackref=true,%
	pdfnewwindow=true,%
	bookmarksnumbered=true,%
	pdfstartview={FitH},
	urlcolor=black,
	linkcolor=black,
	bookmarksopen]%
{hyperref}

\pagestyle{empty}
\setmarginsrb
        {10mm} % left margin
        {10mm} % top margin
        {10mm} % right margin
        {10mm} % bottom margin
        {00mm} % head height
        {00mm} % head sep
        {00mm} % foot height
        {00mm} % foot skip

% Refine sections:
\let\mySection\section
    \renewcommand{\section}[1]{
    \addcontentsline{toc}{section}{#1}%
    \mySection*{#1 \hrulefill}%
}

% Redefine subsections:
\let\mySubsection\subsection
    \renewcommand{\subsection}[1]{%
    \addcontentsline{toc}{subsection}{#1}%
    \mySubsection*{\hspace{5mm}› #1}%
}

% Redefine items:
\let\myItem\item
    \renewcommand{\item}[1]{%
    \myItem {\titlefont #1}%
}

% Custom commands:
\newcommand\email{\begingroup\Url}
\newcommand{\website}[1]{\ \\\hspace{35mm}\email{#1}}
\newcommand{\job}[6]{%
    \item{#1}, #2 (#3), #4, #5%
    \begin{quote}\hspace{-7mm}#6\end{quote}%
}

\begin{document}

\sf

\begin{center}
	{\titlefont\Huge Curriculum Vitæ • Nicolas Dubois}\\
	{\titlefont\hspace{2cm}@wo0dyn, 38 ans, développeur web full stack\footnote{\sf
        Les références pour mes différents postes ou les dissertations
        sont disponibles sur demande.}}
\end{center}

\section{État civil}
    \begin{minipage}[c]{32mm}
        \includegraphics[width=30mm]{images/wo0dyn-circle.jpg}%
    \end{minipage}\hfill%
    \begin{minipage}[c]{20cm}
        \begin{tabular}{p{48mm}@{ }p{80mm}}
            {\bfseries Situation de famille}\dotfill&Divorcé, 2 enfants\\
            {\bfseries Nationalité}\dotfill&Française\\
            {\bfseries Date de naissance}\dotfill&15 juin 1979 à Nancy (54)\\
            {\bfseries Courriel}\dotfill&\email{nicolas.c.dubois@gmail.com}\\
            {\bfseries Page personnelle}\dotfill&\email{https://nicolasdubois.com}\\
            {\bfseries Moyens de locomotion}\dotfill&Permis B, véhicule\\
            {\bfseries Centres d'intérêts}\dotfill&Musique (guitare),
                sports (basket-ball, golf, ski)\\
        \end{tabular}
    \end{minipage}%

\section{Expérience professionnelle}

\subsection{Emplois en informatique}

\begin{description}

    \job{Oscaro}{télétravail, Gennevilliers}{91}{CDI}{Avril 2014--Novembre 2015}{
        Développement du site international d'Oscaro sur une stack Linux,
        Python (Django), PostgreSQL, Redis, ElasticSearch.
        \website{https://www.oscaro.be}
    }

    \job{AMG Développement | Groupe GPdis France}{Eurocentre}{31}{CDI}{Octobre 2009--Avril 2014}{
        J'ai été développeur principal d'applications web chez AMG Développement,
        pôle e-commerce de la filiale du groupe GPdis France.
        Je suis intervenu sur les applications internes (Django, symfony, ZF)
        et également sur la plateforme Magento du groupe (Discounteo, Villatech, Pulsat).
        \website{http://www.villatech.fr}
    }

    \job{Ekinos | Groupe Eklas}{Labège}{31}{CDI}{Avril 2009--Octobre 2009}{
        J'ai été développeur/concepteur web et chef de projet junior pour la société Ekinos,
        une agence web spécialisée dans le développement de solutions e-commerce avec la
        plateforme Magento.
        \website{http://www.labsoft.fr}
    }

    \job{WaterProof}{Labastide Saint-Pierre}{82}{CDI}{Novembre 2007--Avril 2009}{
        J'ai été développeur/concepteur web pour la société WaterProof, où j'ai
        participé au développement d'applications web à forte valeur ajoutée pour
        des grands comptes (Spot Image, Dalta, ITM) de la région toulousaine avec
        le framework symfony\footnote{\sf J'ai suivi la formation symfony
        dispensée par SensioLabs, créateur du framework.}.
        \website{http://www.waterproof.fr}
    }

    \job{MaisMoinsCher.com}{Gaillac}{81}{CDI}{Novembre 2005--Octobre 2007}{
        J'ai été développeur web principal pour le e-commerçant MaisMoinsCher.com
        où j'ai développé :\\
            \hspace*{5mm}- une application intranet pour la gestion du catalogue
                de la boutique ;\\
            \hspace*{5mm}- plusieurs modules dans la partie backoffice dont
                notamment un gestionnaire de stock ;\\
            \hspace*{5mm}- un module de paiement Sofinco pour la partie
                frontoffice ;\\
            \hspace*{5mm}- un pricebot (outil de veille concurrentielle) pour la
                gestion des fournisseurs.
        \website{http://www.maismoinscher.com}
    }

    \job{Les cahiers de la cypriolette}{Commecy}{55}{CDD}{Août 2005}{
        Développement de scripts pour interfacer une base de données MySQL avec
        une application Flash.
        \website{http://www.cypriolette.com}
    }

    \job{Laboratoire Leibniz--IMAG}{Grenoble}{38}{Stage Master de Recherche}{Février--Juillet 2005}{
        Développement et intégration d'un module émotionnel sur une plateforme
        d'agent conversationnel de recommandation.
        Modélisation en UML et implémentation en langage Java.
        \website{http://magma.imag.fr}
    }

\end{description}

\subsection{Conférences}

\begin{description}
        \job{SudWeb (1\iere{} édition)}{Nîmes}{30}{LT}{Mai 2011}{
            Vers des langues de programmation.
            (\email{http://sudweb.fr/2011/})
        }

        \job{Pytong}{Toulon}{83}{LT}{Juin 2013}{
            À quoi ressemblerait mon python ?
            (\email{http://2013.pytong.org})
        }

        \job{DjangoCong}{Clermont-Ferrant}{63}{Conférence}{Juin 2015}{
            Surveiller son application Django avec Sentry.
            (\email{http://rencontres.django-fr.org/2015/})
        }

        \job{DjangoCong}{Clermont-Ferrant}{63}{LT}{Juin 2015}{
            Invoke • Exécuter ses tâches en python.
            (\email{http://rencontres.django-fr.org/2015/})
        }
\end{description}


\subsection{Enseignement}

\begin{description}
        \job{Vidéoscop de Nancy}{Nancy}{54}{Vacation}{Mars--Juin 2003}{
            Médiatisation du cours d'Interface Homme-Machine de Kamel
            Smaïli (PR - Nancy 2) pour le projet e-miage.
            Réalisation de supports multimédia et d'expériences à l'aide des
            outils Flash et PHP.
            \website{http://www.e-miage.org}
        }

        \job{Université Nancy 2}{Nancy}{54}{Tuteur}{Février--Juin 2002}{
            Tuteur en informatique à la faculté d'AES : initiation aux outils de
            bureautique.
            \website{http://www.univ-nancy2.fr/VIDEOSCOP}
        }
\end{description}

\section{Formation}

\subsection{Diplômes}

\begin{description}
    \item{2004--2005 : Master de Recherche « Information, Cognition et Apprentissages »}, %
        spécialité Sciences Cognitives obtenu à l'Institut National Polytechnique
        de Grenoble, mention~AB
        \begin{quote}
            \textbf{Mémoire}\footnote{\sf Mémoire publié aux Éditions
                universitaires européennes, ISBN 978-613-1-59576-9.} :%
                \emph{modélisation informatique de la prise de décision à base d'émotions.}
        \end{quote}
    \item{2003--2004 : Maîtrise de Sciences Cognitives} obtenue à l'Université Nancy 2, mention~AB
        \begin{quote}
        \textbf{Travail d'Études et de Recherche} :
            \emph{état de l'art des systèmes de communication assistée par ordinateur,
                point de vue ergonomie cognitive et ingénierie des langues}
        \end{quote}
    \item{2002--2003 : Licence de Sciences Cognitives} obtenue à l'Université Nancy 2
    \item{2000--2001 : DEUG MIAGE} obtenu à l'IUP MIAGE\footnote{\sf IUP MIAGE :
        Institut Universitaire Professionnalisé en Méthodes Informatiques Appliquées à la
        Gestion des Entreprises.} de l'Université Nancy 2
    \item{1998--1999 : Baccalauréat Scientifique} obtenu au Lycée Henri Poincaré de Nancy
\end{description}

\section{Compétences}

\subsection{Compétences informatiques}

\begin{description}
    \item{Environnements : }\textbf{GNU/Linux} (Ubuntu, Debian), Unix (Solaris, SunOS), MacOS X ;
    \item{Langages : }\textbf{Python}, Unix Shells (bash/zsh), Java, {\LaTeX}, XML, XSL ;
    \item{Développement Web : }\textbf{Python}, \textbf{HTML} 5, (S)\textbf{CSS}, JavaScript, PHP ;
    \item{Bases de Données : }\textbf{PostgreSQL}, MySQL, SQLite ;
    \item{Applications/Frameworks : }\textbf{Django}, flask, symfony (1.x), Magento.
    \item{Notions : }ElasticSearch, Redis, React.js.
\end{description}

\subsection{Compétences linguistiques}

\textbf{Français : }langue maternelle • \textbf{Anglais : }lu, écrit, parlé • \textbf{Espagnol : }notions.

\end{document}
